% !TeX root = ../main.tex
\section{Amplifier and Effect Modelling}
\label{sec:modelling}

Modelling is the forward problem: given the dry signal, produce the processed
one. It matters here for two reasons. A model is one possible output of a tone
estimator (a capture loaded into a player), and a differentiable or
low-dimensional model makes the inverse problem easier to search.

\subsection{Black-box models from paired data}

Feed-forward WaveNets~\cite{damskagg2019deep,damskagg2019realtime} and
recurrent networks~\cite{wright2019realtime} model amplifiers and pedals from
a few minutes of paired audio. Wright et al.\ report test ESR of 0.2--0.3\%
and MUSHRA scores of 90--98 on two tube amplifiers, faster than real time on a
CPU~\cite{wright2020realtime}; a perceptually weighted pre-emphasis improves
the result at no cost~\cite{wright2020perceptual}. Temporal convolutional
networks with fast-growing dilations model compressors~\cite{steinmetz2022efficient},
and time-varying feature modulation extends the memory of such models for fuzz
and compression~\cite{comunita2023modelling}. The largest comparison, 40
devices with LSTM, TCN, GCN and S4 models, ranks state-space models with
time-varying modulation best and finds LSTMs unreliable across
devices~\cite{comunita2025differentiable}. Vanhatalo et al.\ review the
field~\cite{vanhatalo2022review}.

Neural Amp Modeler (NAM) has become the open format for such captures: an
MIT-licensed trainer, core and plugin~\cite{atkinsonndnam}, a JSON
\texttt{.nam} file with WaveNet (``A1'', and since 2026 the slimmable
``A2''~\cite{atkinson2025slimmable,atkinson2026a2}) and LSTM architectures,
and optional metadata such as gear type, tone type (clean, crunch, high gain,
fuzz) and gain. The metadata are ready-made weak labels for a capture library.
The Rust runtime nam-rs runs a standard WaveNet at about 11 times real time on
one x86 core and does not resample, so models must match the host's rate. Any
capture is a snapshot of one setting of one device, made with the device
present.

\subsection{Controls and many devices in one model}

A model can take the knobs as inputs. Juvela et al.\ recorded 4.5\,h of an
amplifier with a robot turning its knobs and trained a conditioned LSTM that
listeners could not tell from an offline circuit simulation~\cite{juvela2023endtoend}.
PANAMA chooses which settings to record by the disagreement of an ensemble,
and with 75 settings matches per-setting NAM captures in a listening
test~\cite{grotschla2025parametric}. Hypernetwork conditioning beats
concatenation for recurrent models~\cite{yeh2024hyper}.

One model can also cover many devices. Chen et al.\ condition a GCN on a
tone embedding from a contrastively trained encoder, so that a wet reference
clip selects the tone; on unseen amplifiers the loss is about twice the
seen-amplifier loss, and a case study missed an equaliser that the model did
not represent~\cite{chen2024zeroshot,chen2024demo}. Open-Amp trains a TCN with
a 394-entry device embedding on renders through crowd-sourced captures; an
unseen analog pedal is added by fitting only its embedding, which matches
one-to-one training with little data~\cite{wright2025openamp}. Embedding-only
enrolment is a small search space, which makes it a candidate target for a
wet-only loss, although the paper used paired data.

\subsection{Grey-box and differentiable signal processing}

Grey-box models keep the structure of the circuit. Block-oriented
Wiener--Hammerstein models (linear filter, static nonlinearity, linear filter)
were identified from measurements~\cite{eichas2017block,eichas2018graybox},
then made differentiable~\cite{kuznetsov2020differentiable,colonel2022waveshapers}.
Esqueda et al.\ learn the component values of the Fender, Marshall and Vox
tone stack from audio~\cite{esqueda2021differentiable}. Miklánek et al.\ put a
white-box tone stack between a recurrent preamplifier and power stage; trained
on 4--12 minutes covering one to three tone-stack settings, it scored 87--95
MUSHRA on unseen settings where a black-box network trained on the same data
scored 18--47~\cite{miklanek2023neural}. The DDSP guitar amplifier of Yeh et
al.\ (preamplifier, tone stack, push--pull power amplifier, output
transformer) needs about 1,350 operations per sample, beats a recurrent model
of the same cost, and trails a model 15 times more
expensive~\cite{yeh2025ddsp}. Grey-box compressors~\cite{wright2022greybox},
phasers~\cite{carson2023differentiable} and modulation
effects~\cite{carson2026gradient} follow the same idea. The comparison of
Comunità et al.\ finds grey-box models not yet as accurate as neural networks
for distortion and fuzz, though competitive for
compressors~\cite{comunita2025differentiable}.

Libraries make these models reusable: DDSP~\cite{engel2020ddsp},
dasp-pytorch and DeepAFx-ST~\cite{steinmetz2022style}, GRAFX~\cite{lee2024grafx},
NablAFx~\cite{comunita2025nablafx} and DiffVox~\cite{yu2025diffvox}; Hayes et
al.\ review differentiable signal processing~\cite{hayes2024review}. DiffVox
fits a differentiable equaliser, compressor, delay and reverberator to each of
435 vocal tracks and studies the distribution of the fitted presets, a recipe
that would turn guitar stems into a prior over settings.

The structure of these models resembles the amplifier of a typical effects
application (gain, voicing, three-band tone stack). That matters less for
modelling quality than for estimation: a few interpretable parameters are
what a search can explore and a player can adjust.
